\section{The full method family}\label{sec:family}

This section extends the analysis to the remaining methods of \cite{graf2026numerical}:
hidden-variable resultants (with \cite{noferini2016numerical}), Gr\"obner elimination, the
rational univariate representation, multiparameter eigenvalue problems, and the Macaulay
resultant. The pipeline of Sections~\ref{sec:jacobian}--\ref{sec:local} computes the local
terms of any method; the methods differ in their structural terms, which we now classify.

\subsection{Hidden-variable resultants are local methods}

The key input is an identity proved in \cite{noferini2016numerical} which, read with the
present machinery, moves the hidden-variable resultant methods from the pairwise class to
the local one.

\begin{lemma}[\cite{noferini2016numerical}]\label{lem:hiddenvariable}
  Let $R(x_d)$ be the Cayley resultant matrix polynomial of the system in the monomial
  basis, for any $d\ge2$, or the Sylvester matrix polynomial for $d=2$. Let $z_j$ be a
  simple root whose $d$-th coordinate $x_{d,j}$ is a simple eigenvalue of $R$, and let
  $v_j,w_j$ be the eigenvectors in the canonical normalizations of
  \cite[Lem.~3.5]{noferini2016numerical} (Vandermonde, for Cayley) and
  \cite[Lem.~4.4]{noferini2016numerical} (Vandermonde and Horner shifts, for Sylvester).
  Then
  \[
    w_j^T\,R'(x_{d,j})\,v_j = \det J(z_j)
  \]
  exactly, and therefore, by the eigenvalue condition number formula
  \cite[Lem.~2.11]{noferini2016numerical},
  \[
    \kappa(x_{d,j},R) = \frac{\|v_j\|_2\,\|w_j\|_2}{|\det J(z_j)|},
    \qquad
    \frac{\kappa(x_{d,j},R)}{\kroot(z_j)}
    = \frac{\sigma_{\min}(J(z_j))\,\|v_j\|_2\,\|w_j\|_2}{|\det J(z_j)|}.
  \]
\end{lemma}

\proof
The identity is \cite[Thm.~3.6]{noferini2016numerical} in matrix form for Cayley, and is
established inside the proof of \cite[Lem.~4.5]{noferini2016numerical} for Sylvester. The
condition number formula is an equality, and the ratio $\|v\|\|w\|/|w^TR'v|$ is invariant
under rescaling of the eigenvectors, so any normalization gives the same value.
\endproof

Every factor on the right is a local functional of the root: $\sigma_{\min}(J)$ and
$\det J$ by Lemma~\ref{lem:jacobian}, $\|v_j\|$ and $\|w_j\|$ because the canonical
eigenvectors are built from the root coordinates (Cayley) or from the fiber polynomials at
the root (Sylvester). No pair energy and no deflated determinant survive; comparing with
Lemma~\ref{lem:adjugate}(a) below gives the byproduct identity
$\gamma_j = r'(x_j)/\det J(z_j)$ in the canonical normalization, so the two exponentially
large quantities $r'(x_j)$ and $\gamma_j$ cancel exactly. The scalar-eliminant methods of
Section~\ref{sec:elimination} do retain a pair energy; the matrix formulation is precisely
what removes it, a quantitative vindication of the ``numerically confirmed eliminants''
terminology recalled in \cite[Sec.~2.2]{noferini2016numerical}.

\subsection{The Cayley average: an unconditional theorem}\label{sec:cayley}

In the model of Definition~\ref{def:model} the Cayley construction is generically regular:
roots at infinity and singular resultant matrix polynomials form a null set, by the box
Newton polytopes. For Cayley in the monomial basis,
\cite[Lem.~3.5]{noferini2016numerical} gives $\|v_j\|$ and $\|w_j\|$ in closed form: with
$\tau_k = kn-1$,
\[
  \|v_j\|_2^2 = \prod_{k=1}^{d-1}k_{\tau_{d-k}}(r_k), \qquad
  \|w_j\|_2^2 = \prod_{k=1}^{d-1}k_{\tau_k}(r_k), \qquad
  r_k = |(z_j)_k|^2.
\]

\begin{theorem}\label{thm:cayley}
  In the model of Definition~\ref{def:model}, for fixed $d\ge2$ and $n\to\infty$, with
  $\kappa(\cdot,R_{\mathrm{Cay}})$ the absolute condition number of
  \cite{noferini2016numerical} in the monomial basis,
  \[
    \frac1N\,\E\sum_{j=1}^N\log\frac{\kappa(x_{d,j},R_{\mathrm{Cay}})}{\kroot(z_j)}
    \;=\; -\,\frac{d(d-1)}2\,\log n \;+\; O_d(1).
  \]
\end{theorem}

\proof
By Lemma~\ref{lem:hiddenvariable} the summand is
$\log\sigma_{\min}(J) + \log\|v_j\| + \log\|w_j\| - \log|\det J|$, a local functional, so
Lemma~\ref{lem:rice} applies; the integrability checks of Lemma~\ref{lem:integrability}
extend verbatim to $\log|\det J|$ and $\log k_m(r_l)$. The four pieces are:
\begin{enumerate}[(1)]
  \item $\Eroots[\log\sigma_{\min}(J)] = (d+1)\log n + O(d)$, by
    Proposition~\ref{prop:local}.
  \item $\Eroots[\log|\det J|] = \EW[\log|\det G|]
    + \frac d2\sum_l\Eroots[\log k_n(r_l)] + \frac12\sum_l\Eroots[\log h(r_l)]$, since
    $|\det J| = |\det G|\,(\det\Sigma)^{1/2}$ and $\det\Sigma = S^d\prod_lh(r_l)$; by
    Corollary~\ref{cor:weighted}(a) and Proposition~\ref{prop:outer} this equals
    $d(d+1)\log n + \frac12\sum_{j=2}^{d+1}\psi(j) + O(d^2)$.
  \item $\Eroots[\log\|v_j\| + \log\|w_j\|]
    = \frac12\sum_{k=1}^{d-1}\bigl(\E[\log k_{(d-k)n-1}(r_k)]
    + \E[\log k_{kn-1}(r_k)]\bigr)
    = \frac12\sum_{k=1}^{d-1}(d+2)\log n + O_d(1)
    = \frac{(d-1)(d+2)}2\log n + O_d(1)$, by Corollary~\ref{cor:logkm} with $a=d-k$ and
    $a=k$.
\end{enumerate}
Adding: $(d+1) - d(d+1) + \frac{(d-1)(d+2)}2 = -(d^2-1) + \frac{d^2+d-2}2
= -\frac{d(d-1)}2$.
\endproof

\begin{remark}\label{rem:cayley}
  Four remarks are in order. First, the sign: on average, in the absolute metric of
  \cite{noferini2016numerical}, the eigenvalue of the Cayley eigenproblem is \emph{less}
  sensitive than the root, by a polynomial factor; the exponential worst case
  $\kappa\ge\kroot^d$ of \cite[Thm.~3.7]{noferini2016numerical} is a tail event in this
  model. Second, the units: absolute condition numbers compare perturbations of $R$ with
  perturbations of $p$, which live at different scales, and
  \cite[Sec.~5]{noferini2016numerical} shows that the relative comparison can behave
  differently; Section~\ref{sec:relativeCayley} carries out the relative analysis and shows
  that the apparent paradox disappears there. Third, the basis normalization
  $\sup_\Omega|\phi_k| = 1$ of \cite{noferini2016numerical} is satisfied by monomials on
  the closed unit polydisk, which contains the bulk of the roots; the outer roots are
  handled by the same tail estimates as everywhere above. Fourth, the theorem concerns the
  eigenvalue subproblem only: the construction of $R$ and the eigenvector extraction are
  outside its scope, exactly as in
  \cite{noferini2016numerical,graf2026numerical}.
\end{remark}

For Sylvester at $d=2$ the same reduction gives
\[
  \frac1N\,\E\sum_j\log\frac{\kappa(x_{2,j},R_{\mathrm{Sylv}})}{\kroot(z_j)}
  = 3\log n - \bigl(6\log n + O(1)\bigr) + \tfrac32\log n
  + \Eroots\log\|w_j\| + O(1),
\]
by Proposition~\ref{prop:local}, item (2) of the proof above, and
Corollary~\ref{cor:logkm} with $a=2$ for $\|v\| = \|(x_1^{2n-1},\dots,1)\|$. The remaining
factor is the Horner-shift norm, and here, uniquely in this paper, the conditioning on the
root matters at leading order. The entries of $w_j$ are the coefficients of the deflated
fibers, $b_k = \sum_{j>k}a_jx_1^{j-k-1}$ with fiber coefficients $a_j$ of variance
$k_n(r_2)$, and the root constraint $\sum_ja_jx_1^j = 0$ is one linear condition on them;
the rank-one Schur complement gives, exactly,
\[
  \E_{\mathrm{cond}}\|w_j\|^2 = 2\,k_n(r_2)\,G_c(r_1), \qquad
  G_c(r) = \frac{n\,(r^{n+1}+1) - 2r\,\dfrac{r^n-1}{r-1}}{(r-1)\,(r^{n+1}-1)}, \qquad
  G_c(1) = \frac{n(n+2)}6.
\]
On the window $r=e^{s/n}$ one has $G_c = n^2\gamma_c(s)(1+o(1))$ with
$\gamma_c(s) = \bigl[(e^s+1) - 2(e^s-1)/s\bigr]/\bigl[s(e^s-1)\bigr]$, and
$\gamma_c(s)\sim1/s$ as $s\to\infty$: the constraint cancels the exponential growth of the
deflated coefficients at outer roots, so, in contrast with every unconditioned monomial
functional above, there is \emph{no} outer inflation. Passing from the second moment to the
log-average costs only a universal constant, with no assumption on the conditional
spectrum: diagonalizing the conditional covariance and writing
$\|w_j\|^2 = \sum_i\lambda_i|g_i|^2$ with independent standard complex Gaussians $g_i$, the
concavity of the logarithm gives pointwise
$\log\|w_j\|^2\ge\log\sum_i\lambda_i + \sum_i\frac{\lambda_i}{\sum\lambda}\log|g_i|^2$, so
that
$\log\E_{\mathrm{cond}}\|w_j\|^2 - \gamma
\le\E_{\mathrm{cond}}\log\|w_j\|^2\le\log\E_{\mathrm{cond}}\|w_j\|^2$, using
$\E\log|g|^2 = -\gamma$. Hence $\Eroots\log\|w_j\| = 2\log n + O(1)$ (by quadrature,
$\E\log G_c - 2\log n\to-2.25$), and the Sylvester coefficient is
\[
  3 - 6 + \tfrac32 + 2 = +\tfrac12:
\]
on average, in the absolute metric, the Sylvester eigenvalue is worse conditioned than the
root by $\sqrt n$ and worse than the Cayley eigenvalue by $n^{3/2}$, consistent in
direction with the $\|v\|/\|w\|$ comparison of \cite[Sec.~4.3]{noferini2016numerical}. An
unconditioned computation would give $+1$; the difference is real and visible in the data.
The one ingredient not captured by the Schur complement is the $|\det J|^2$ weight of the
root average; we rely here on the numerically supported assumption that its effect is a
bounded shift, measured at $+0.08$ to $+0.12$: the finite-$n$ prediction built from $G_c$
tracks the measured $\E\log\|w_j\| = 2.21, 2.78, 3.24, 3.80$ at $n=4,6,8,11$ within that
margin, and the assembled ratio prediction tracks the measured $0.73, 0.98, 1.19, 1.38$
within $0.06$--$0.12$.

\subsection{Numerical verification of Theorem \ref{thm:cayley}}\label{sec:verification}

All experiments were run in double precision in Python (NumPy/SciPy, LAPACK backends):
polynomial eigenvalue problems are solved by companion linearization and the QZ algorithm;
roots are refined by Newton iteration to residuals near machine precision;
Vandermonde-type systems are assembled column-normalized in log-space, which is legitimate
since $\Lambda_j$ is invariant under column scaling, and solved by LU factorization with
explicit residual monitoring; deterministic predictions are evaluated by adaptive
quadrature, with extended-precision arithmetic near removable singularities. The scripts
accompany this paper.

Because the sign of Theorem~\ref{thm:cayley} is surprising, we tested it directly at
$d=2$: sample the model, build the B\'ezout (Cayley) matrix polynomial $B(x_2)$ of size
$n$, compute all roots by linearization plus Newton refinement, and evaluate every quantity
per root. Three outcomes, for $n\in\{4,6,8,11,14\}$ with $20$--$60$ samples per $n$ (all
$2n^2$ roots recovered in every run):
\begin{enumerate}[(1)]
  \item The identity of Lemma~\ref{lem:hiddenvariable} holds to machine precision: the
    eigenvalue condition number computed from numerical null vectors of $B(x_2^*)$ and the
    Rayleigh quotient agrees with $\|v\|\|w\|/|\det J|$ with relative error below
    $10^{-11}$ at every root.
  \item The measured averages match the \emph{finite-$n$} deterministic predictions of the
    theory (the integrals $\frac1n\int h_n\log k_m\,dr$ etc., evaluated by quadrature, plus
    the tilted Gaussian constants); see Table~\ref{tab:cayley}. The residual gap in the
    last column is the min--max slack of the $\sigma_{\min}$ sandwich \eqref{eq:sandwich},
    which the single-coordinate prediction does not resolve; it is bounded and of the
    observed size.
  \item Fitting a slope in $\log n$ over $n\le14$ gives $-0.58$ rather than $-1$; this is
    a finite-size effect, not a discrepancy.
  \item The same deterministic prediction, evaluated for large $n$, converges to
    $-\log n + 1.12 + o(1)$: the constant inside the theory converges only logarithmically
    (its per-coordinate part is now closed-form, Proposition~\ref{prop:closedform}), and
    at $n=14$ about $40\%$ of it is still missing. The theorem's coefficient
    $-\frac{d(d-1)}2 = -1$ is thereby confirmed by the $n\to\infty$ behaviour of the exact
    integrals, while the finite-$n$ data sit on the exact finite-$n$ curve.
\end{enumerate}

\begin{table}[ht]
  \centering
  \begin{tabular}{r|cc|cc|cc}
    & \multicolumn{2}{c|}{$\Eroots\log|\det J|$}
    & \multicolumn{2}{c|}{$\Eroots\log\|v\|\|w\|$}
    & \multicolumn{2}{c}{$\Eroots\log(\kappa_{\mathrm{Cay}}/\kroot)$} \\
    $n$ & pred. & meas. & pred. & meas. & pred. & meas. \\
    \hline
    4 & 5.39 & 5.39 & 1.93 & 1.90 & $-1.18$ & $-1.41$ \\
    8 & 8.41 & 8.42 & 3.06 & 3.11 & $-1.55$ & $-1.67$ \\
    14 & 11.11 & 11.11 & 4.04 & 3.99 & $-1.93$ & $-2.16$
  \end{tabular}
  \caption{Measured root averages at $d=2$ for the Cayley method against the finite-$n$
    deterministic predictions of the theory.}
  \label{tab:cayley}
\end{table}

Five further measurements were carried out with the same infrastructure.

\emph{Cayley at $d=3$.} The trivariate Cayley tensor was built by convolution and two
B\'ezout-type divisions, and self-tested: the identity of Lemma~\ref{lem:hiddenvariable}
holds to $10^{-11}$ in the log across all roots, and all $6n^3$ roots are recovered at
$n=2,3,4$, validating the construction. The measured ratio
$\Eroots\log(\kappa_{\mathrm{Cay}}/\kroot)$ is $-2.53, -3.00, -3.29$, negative and
monotonically decreasing, consistent in sign and direction with the coefficient
$-\frac{d(d-1)}2 = -3$ of Theorem~\ref{thm:cayley}, with the strong finite-size
attenuation expected at such small $n$.

\emph{Sylvester.} The identity of Lemma~\ref{lem:hiddenvariable} for the Sylvester matrix
holds to relative error below $4\cdot10^{-3}$ in the log (degrading with the conditioning
of $S(x_2^*)$ at outer roots, but negligible against the means), and the measured
$\Eroots\log(\kappa_{\mathrm{Sylv}}/\kroot)$ is positive and increasing,
$0.73, 0.98, 1.19, 1.38$ at $n=4,6,8,11$, and is reproduced within $0.06$--$0.12$ by the
finite-$n$ prediction built from the conditional fiber moment $G_c$ of
Section~\ref{sec:cayley}, whose asymptotic coefficient is $+\frac12$.

\emph{Coordinate pair energy.} The coordinate pair energy
$\frac1N\sum_j\sum_{i\ne j}\log|x_{2,j}-x_{2,i}|$, the structural term of Gr\"obner
elimination, measures $9.9, 17.4, 29.0, 46.2$ at the same $n$, against the
independent-pair quadrature prediction $9.3, 17.4, 26.6, 42.2$: correlations contribute
under ten percent, and the growth is $\Theta(n\log n)$, an outer-root effect again, since
$\E\log|x-x'|$ for an independent pair equals
$\E\log\max(|x|,|x'|)\approx(\log n)/n$ per pair, times $N-1 = 2n^2-1$ pairs.

\emph{Projected pair energy.} The projected pair energy of the rational univariate
representation with balanced $u$: measured $11.1, 19.9, 31.4, 50.4$ at $n=4,6,8,11$,
against the one-point Monte Carlo $(N-1)I_n = 10.5, 19.6, 29.1, 48.1$, with
$I_n\to I(u) = 7\zeta(3)/(2\pi^2) - \frac12\log2$ (Proposition~\ref{prop:rur});
correlations again contribute under five percent.

\emph{Overlap term.} The overlap term $\Theta_n$ at $d=1$ and $d=2$; this is reported in
Section~\ref{sec:protocol}, where it decides between the two candidate regimes. The
scripts are short (a B\'ezout builder, a companion linearization, Newton refinement, and
quadrature routines) and are available alongside this paper.

\subsection{Relative conditioning: resolution of the apparent paradox}
\label{sec:relativeCayley}

Theorem~\ref{thm:cayley} says that in the absolute metric the Cayley eigenvalue is
typically \emph{better} conditioned than the root. This is not absurd, but it is a
statement about units: an absolute perturbation of size $\varepsilon$ means something
different for the resultant matrix, whose entries are degree-$d$ forms in the coefficients
and are large, than for the polynomials themselves. Following
\cite[Sec.~5]{noferini2016numerical}, the relative condition numbers are
$\kappa_{\mathrm{rel}} = \kappa_{\mathrm{abs}}\cdot\|D\|/\|X\|$ with $D$ the data and $X$
the solution, so
\begin{multline*}
  \Eroots\log\frac{\kappa_{\mathrm{rel}}(x_d,R_{\mathrm{Cay}})}{\kappa_{\mathrm{rel,root}}}
  = \Eroots\log\frac{\kappa(x_d,R_{\mathrm{Cay}})}{\kroot} \\
  + \E\log\sup_{|x_d|\le1}\|R_{\mathrm{Cay}}(x_d)\|_2
  - \E\log\max_i\sup_\Omega|p_i| + O(1),
\end{multline*}
the $O(1)$ collecting $\E\log(\|x^*\|/|x_d^*|)$.

The data scales are Gaussian quantities. For the polynomials,
$\E\log\sup_\Omega|p_i| = \frac d2\log n + O(\log(dn))$, since the field variance on the
torus is $(n+1)^d$. For the resultant matrix, the Frobenius norm of the coefficient tensor
$A$ carries a near-diagonal enhancement that a naive Parseval count misses: at $d=2$,
Parseval in all three variables and the exact two-point formula
\[
  \E|f_{\mathrm{Cay}}(s,t,x)|^2
  = \frac{2\,(n+1)^2\bigl[(n+1)^2 - |k_n(s\bar t)|^2\bigr]}{|s-t|^2}
\]
on the torus give $\E\|A\|_F^2\asymp n^5$: the Dirichlet-kernel integral over $s\bar t$
contributes a factor $n^3$, not $n^2$. In general each eliminated pair contributes one
extra factor of $n$, so $\E\|A\|_F^2\asymp_dn^{d^2+d-1}$.

The operator norm behaves differently, and this is the key point, which can be closed for
every $d$.

\begin{proposition}\label{prop:opnorm}
  For every fixed $d$,
  $\E\log\sup_{|x|=1}\|R_{\mathrm{Cay}}(x)\|_2 = \frac{d^2}2\log n + O_d(\log\log n)$, and
  the coefficient of $\log n$ in the average relative Cayley ratio is $0$.
\end{proposition}

\proofsketch
Upper bound, by induction on $d$. Expanding the Cayley determinant along its first row
gives $f_d = \sum_j\Delta_{s_1,t_1}(p_j)\cdot g_j$, where $\Delta_{s_1,t_1}$ is the divided
difference in the first eliminated pair and $g_j$ is the Cayley function of the $j$-th
$(d-1)\times(d-1)$ minor, which is divisible by each remaining $(s_i-t_i)$ individually
because its defining rows coincide there. Hence
$R_d(x_d) = \sum_jH(\Delta p_j)\circ R_{d-1}^{(j)}$, a sum of compositions of triangular
Hankel--Toeplitz sections whose symbols are the polynomials on subtori and lower-order
Cayley matrices; at $d=2$ this is the classical Gohberg--Semencul-type factorization of the
B\'ezout matrix. Each finite triangular section of a bounded symbol costs a factor
$C\log n$ (finite-section Riesz projection), and by Salem--Zygmund with Gaussian
concentration $\E\log\sup_{\mathbb T^d}|p_j| = \frac d2\log n + O(\log\log n)$; unwinding
the induction,
$\log\sup_x\|R_d(x)\|_2\le\sum_{j=1}^d\log\sup|p_j| + O_d(\log\log n)
= \frac{d^2}2\log n + O_d(\log\log n)$ in expectation. Lower bound:
$\|R(x)\|_2\ge\|R(x)\|_F/\sqrt{(d-1)!\,n^{d-1}}$, with
$\E\|R(x)\|_F^2\asymp_dn^{d^2+d-1}$ (exact at $d=2$ by the two-point formula above, one
near-diagonal factor $n$ per eliminated pair in general, at the sketch level already
flagged), and the lower deviation of the degree-$d$ Gaussian chaos $\|R(x)\|_F^2$ is
controlled by Carbery--Wright, so
$\E\log\|R(x)\|_2\ge\frac{d^2}2\log n - O_d(1)$ at any fixed torus point.
\endproof

The spectrum of $R(x)$ on the torus is flat: with rank $(d-1)!\,n^{d-1}$, the
flat-spectrum operator norm is
\[
  n^{\frac{d^2+d-1}2 - \frac{d-1}2} = n^{\frac{d^2}2},
\]
the near-diagonal enhancement of the Frobenius norm exactly compensating the rank
normalization. The measured exponents at $d=2$ confirm both statements: the fitted slope
of $\log\sup_{|x|=1}\|R(x)\|_2$ over $n\le24$ is $2.05$ against $d^2/2 = 2$, while
$\log\|R(x)\|_F$ and $\log\|A\|_F$ fit $2.4$ and drift toward $(d^2+1)/2 = 2.5$. A
rigorous lower bound consistent with this comes from testing against the monomial vectors
$v_s,v_t$ of the two groups of Cayley variables, for which
$f_{\mathrm{Cay}}(s,t,x) = v_s^TR(x)\,v_t$; the matching upper bound is the content of
Proposition~\ref{prop:opnorm}.

Taking $\E\log\sup_x\|R(x)\|_2 = \frac{d^2}2\log n + O_d(1)$, the coefficient of $\log n$
in the average relative ratio is
\[
  -\frac{d(d-1)}2 + \frac{d^2}2 - \frac d2 = 0:
\]
\emph{relative neutrality}. This is now supported on two independent fronts: the direct
measurement of the relative ratio in Section~\ref{sec:verification}, which increases
toward and then settles near $0.6$ in natural-log units over the accessible range, and the
measured operator-norm exponent above. The picture is coherent: absolutely, the Cayley
eigenvalue is typically better conditioned than the root by the polynomial factor of
Theorem~\ref{thm:cayley}; relatively, the two are typically comparable; and the
exponential worst cases of \cite{noferini2016numerical} are tail events in both metrics.

\subsection{The general factorization, for other resultant matrices}

Lemma~\ref{lem:hiddenvariable} is specific to the Cayley and Sylvester matrices with their
canonical eigenvector normalizations. For other resultant matrices, or other weightings,
the following general factorization applies; it is superseded by
Lemma~\ref{lem:hiddenvariable} in the two canonical cases, through the identity
$\gamma_j = r'(x_j)/\det J(z_j)$ noted above. Let $S(x)$ be the $2n\times2n$ Sylvester
matrix polynomial in the $y$-coefficients of $p_1(x,\cdot),p_2(x,\cdot)$;
$r(x) := \det S(x) = \Res_y(p_1,p_2)$ has degree $N=2n^2$ and its roots are the
$x$-coordinates of the roots $z_j = (x_j,y_j)$. The method computes the $x_j$ as
eigenvalues of the polynomial eigenvalue problem $S(x)v = 0$. With weights
$\omega_0,\dots,\omega_n$ on the matrix coefficients, the normwise unstructured condition
number of the eigenvalue $x_j$ is \cite{tisseur2000backward}
\begin{equation}\label{eq:pep}
  \kappa_{\mathrm{res}}(x_j)
  = \frac{\bigl(\sum_{k=0}^n\omega_k|x_j|^k\bigr)\,\|w_j\|_2\,\|v_j\|_2}
  {|w_j^*\,S'(x_j)\,v_j|},
\end{equation}
with $v_j,w_j$ right and left null vectors of $S(x_j)$. Unstructured is the honest choice:
the QZ iteration perturbs the matrices, not the original coefficients, and the
structured--unstructured gap is part of the instability mechanism of
\cite{noferini2016numerical}.

\begin{lemma}\label{lem:adjugate}
  Let $x_j$ be a simple root of $r=\det S$ and $v_j,w_j$ null vectors of $S(x_j)$, in any
  normalization. Then $\adj S(x_j) = \gamma_jv_jw_j^*$ for a unique $\gamma_j\ne0$, and:
  \begin{enumerate}[(a)]
    \item $r'(x_j) = \gamma_j\,w_j^*S'(x_j)v_j$;
    \item $\gamma_j\,(w_j^*v_j) = \prod_{\mu\in\spec S(x_j)\setminus\{0\}}\mu$, the
      deflated determinant of $S(x_j)$;
    \item combining (a) and \eqref{eq:pep},
      \[
        \kappa_{\mathrm{res}}(x_j)
        = \Bigl(\sum_{k=0}^n\omega_k|x_j|^k\Bigr)
        \frac{\|v_j\|\,\|w_j\|\,|\gamma_j|}{|r'(x_j)|},
        \qquad r'(x_j) = \lc(r)\prod_{i\ne j}(x_j-x_i),
      \]
      and the right-hand side is invariant under separate rescaling of $v_j$ and $w_j$.
  \end{enumerate}
\end{lemma}

\proof
Simplicity gives $\rank S(x_j) = 2n-1$, so $\adj S(x_j)$ has rank one; from
$S\cdot\adj S = \adj S\cdot S = 0$ at $x_j$, its column space is $\ker S(x_j)$ and its row
space the left kernel, whence $\adj S(x_j) = \gamma_jv_jw_j^*$. (a) is Jacobi's formula
$r'(x) = \tr(\adj S(x)\,S'(x))$ evaluated at $x_j$. (b): for any matrix $A$ with a simple
eigenvalue $0$ and null vectors $v,w$, one checks on a Schur form that
$\adj A = \bigl(\prod_{\mu\ne0}\mu\bigr)\frac{vw^*}{w^*v}$; comparison gives (b). (c):
substitute (a) into \eqref{eq:pep}; rescaling $v_j\mapsto cv_j$ forces
$\gamma_j\mapsto\gamma_j/c$, leaving $\|v_j\|\|w_j\||\gamma_j|$ unchanged, and similarly
for $w_j$.
\endproof

Each factor of $\kappa_{\mathrm{res}}/\kroot = \sigma_{\min}(J)\,\kappa_{\mathrm{res}}$
sits at a known level:
\begin{enumerate}[(i)]
  \item $\sigma_{\min}(J)$ and the weight $\sum\omega_k|x_j|^k$ are local, computed by
    Proposition~\ref{prop:local} and the argument used for $C_n$.
  \item $\|v_j\|$ is local: for Sylvester, $v_j = (y_j^{2n-1},\dots,y_j,1)^T$, the
    monomial vector of the other coordinate; its log-average is the $d=1$ functional of
    Section~\ref{sec:outer} with degree $2n-1$, outer-root inflation included, hence
    $\Theta(\log n)$ with a constant delivered by the computation behind \eqref{eq:logk}.
  \item $\|w_j\|$ is local on the fiber: $w_j^*S(x_j) = 0$ encodes the cofactor identity
    $u(y)p_1(x_j,y) + t(y)p_2(x_j,y)\equiv0$ with $u\propto p_2(x_j,y)/(y-y_j)$ and
    $t\propto-p_1(x_j,y)/(y-y_j)$. Thus $\|w_j\|$ depends only on the two univariate fiber
    polynomials at $x=x_j$, conditioned to share the root $y_j$; this is a Gaussian
    conditioning on the fiber, within the reach of the same pipeline.
  \item $|r'(x_j)|$ is a pair statistic:
    \[
      \frac1N\,\E\sum_j\log|r'(x_j)|
      = \E\log|\lc(r)| + \frac1N\,\E\sum_j\sum_{i\ne j}\log|x_j-x_i|,
    \]
    the logarithmic pair energy of the $x$-marginal of the root set, computed at leading
    order in Section~\ref{sec:verification}. For the canonical Sylvester and Cayley cases
    this term cancels against $|\gamma_j|$ by Lemma~\ref{lem:hiddenvariable}; it survives
    for the scalar-eliminant methods of Section~\ref{sec:elimination}.
  \item $|\gamma_j|$ is the deflated determinant of $S(x_j)$; in the canonical
    normalization it equals $|r'(x_j)|/|\det J(z_j)|$, so it needs no separate analysis
    there.
\end{enumerate}

\subsection{Gr\"obner elimination and the rational univariate representation}
\label{sec:elimination}

Both methods pass through an intermediate univariate polynomial, and
\cite{graf2026numerical} measures their instability by the conditioning
$\kappa_{\mathrm{uni}} = 1/|g'|$ of its roots.

For Gr\"obner elimination in the variable $x_1$, the eliminant $g$ generates
$I\cap\C[x_1]$; generically $\deg g = N$, its roots are the $x_1$-coordinates, and, for a
generic system, $g$ agrees with $\Res_{x_2,\dots}(p)$ up to a constant. With the monic
normalization $g(x) = \prod_i(x-x_{1,i})$,
\[
  \kappa_{\mathrm{uni}}(x_{1,j})^{-1} = |g'(x_{1,j})| = \prod_{i\ne j}|x_{1,j}-x_{1,i}|,
\]
so
\[
  \frac1N\,\E\sum_j\log\frac{\kappa_{\mathrm{uni}}(x_{1,j})}{\kroot(z_j)}
  = A_n - \frac1N\,\E\sum_j\sum_{i\ne j}\log|x_{1,j}-x_{1,i}|:
\]
the structural term of Gr\"obner elimination is the coordinate pair energy. Note the
contrast with Lemma~\ref{lem:hiddenvariable}: the matrix formulation of elimination is
local, the scalar one is not. The normalization of $g$ matters for
$\kappa_{\mathrm{uni}}$, and monic is the natural theoretical choice, matching the
examples of \cite{graf2026numerical}; other scalings contribute terms of the form
$\E\log|\lc|$, controlled as above. For contrast, the worst case of this method in
\cite[Ex.~3.1]{graf2026numerical} is $\sigma^{-2^d+2}$, while the average is a pair energy
of $N$ points concentrated near a circle, which Section~\ref{sec:verification} shows to be
$2n(\log n+C)(1+o(1))$ at $d=2$: the monic eliminant has exponentially large derivatives
at its roots, so in the absolute metric elimination too is typically \emph{better}
conditioned than the root, one more units effect; the relative version needs the
coefficient scale of the eliminant, a Carbery--Wright quantity.

For the rational univariate representation, one projects by a linear form
$t(x) = u^Tx$ with $\|u\|\le1$ and solves $f(x) = \prod_i(x-t(z_i))$, which is monic, so
\[
  \kappa_{\mathrm{uni}}(t(z_j))^{-1} = \prod_{i\ne j}|t(z_j)-t(z_i)|:
\]
the pair energy of a one-dimensional linear projection of the whole root set. This energy
can be computed at leading order, with a closed-form constant in the balanced case.

\begin{proposition}\label{prop:rur}
  Let $d=2$, $t(x) = u_1x_1+u_2x_2$ with $u_1,u_2>0$, $u_1^2+u_2^2 = 1$. Then
  \[
    \frac1N\,\E\sum_j\sum_{i\ne j}\log|t(z_j)-t(z_i)| = I(u)\,N\,(1+o(1)),
    \qquad I(u) = \E\bigl[\log\max(u_1D,\,u_2D')\bigr],
  \]
  where $D,D'$ are independent copies of $2|\sin(\beta/2)|$ with $\beta$ uniform on
  $(0,2\pi)$, the correlation correction being measured below five percent at the
  accessible sizes, a constant-level bookkeeping step left at the flagged-sketch level. In
  the balanced case $u_1=u_2=1/\sqrt2$,
  \[
    I = \frac{7\,\zeta(3)}{2\pi^2} - \frac{\log2}2 = 0.0797\ldots > 0.
  \]
\end{proposition}

\proof
At leading order the pairs decouple and, by Proposition~\ref{prop:rootscale}, the two
coordinates of a single root become independent with uniform phases on the unit circle.
Writing $e^{i\varphi} - e^{i\varphi'} = e^{i\alpha}D$ with $\alpha$ uniform and
independent of $D = 2|\sin((\varphi-\varphi')/2)|$, the difference of two independent
projected values is $u_1De^{i\alpha} + u_2D'e^{i\alpha'}$. The mean-value property of the
harmonic function $\log|\cdot|$ gives
$\E_\alpha\log|me^{i\alpha}+c| = \log\max(m,|c|)$ for any fixed $c$, and averaging over
one phase yields $I(u)$. In the balanced case, with $F_D(2\sin y) = \frac2\pi\min(y,\pi-y)$,
\[
  \E\log\max(D,D') = \frac8{\pi^2}\int_0^{\pi/2}y\,\log(2\sin y)\,dy
  = \frac{7\,\zeta(3)}{2\pi^2},
\]
by the classical log-sine evaluation
$\int_0^{\pi/2}y\log(2\sin y)\,dy = \frac7{16}\zeta(3)$, and
$\E\log(D/\sqrt2) = -\frac12\log2$ since $\E\log D = 0$.
\endproof

The sign is noteworthy: $I>0$ in the balanced case, so the monic eliminant of the rational
univariate representation has \emph{exponentially large} derivatives at its roots, and in
the absolute metric $\kappa_{\mathrm{uni}}$ is exponentially small on average, the most
extreme units effect in the family. Near-coincidences of projected roots, the mechanism of
the devastating example of \cite[Sec.~4]{graf2026numerical}, are integrated inside $I(u)$
and only shape the constant; the example itself is a tail event. The finite-$n$ version of
the statement, a one-point Monte Carlo of $(N-1)I_n$, matches the measured energies within
five percent (Section~\ref{sec:verification}). For general $u$ the one-dimensional
integral defining $I(u)$ is immediate, and for general $d$ the same harmonic reduction
applies to the outermost phase, leaving a $(d-1)$-fold integral.

\subsection{Multiparameter eigenvalue problems: local, computed for an explicit
  representation}\label{sec:mep}

The multiparameter approach represents each $p_k$ as the determinant of a matrix of
affine-linear forms, $W_k(x) = A_{k,0} + \sum_lx_lA_{k,l}$ with $\det W_k = p_k$, and
solves the resulting multiparameter eigenvalue problem through its operator determinants
\cite{graf2026numerical}. For the representation to be computable directly from the
coefficients one needs a \emph{uniform} determinantal representation in the sense of
\cite{boralevi2017uniform}, with entries affine-linear in the $c_\alpha$ as well; the
minimal size of such a representation is $\Theta(n^{d/2})$ for fixed $d$, and for $d=2$
the explicit \textsc{MinUnif} construction of \cite[Ex.~5.1]{boralevi2017uniform} has size
$2n+1$: $W = \bigl[\begin{smallmatrix}M_x&0\\L&M_y^T\end{smallmatrix}\bigr]$, with
$M_x,M_y$ bidiagonal in $(-x,1)$ and $(-y,1)$, and the coefficients sitting, each exactly
once, in the antitriangular block $L$. The first observation is that the method is local,
for \emph{every} uniform representation, by the adjugate identity.

\begin{proposition}\label{prop:mep}
  Let $z^*$ be a simple root and, for each $k$, let $W_k(z^*)$ have a one-dimensional
  kernel; normalize the null vectors by $\adj(W_k(z^*)) = v_kw_k^T$. Then:
  \begin{enumerate}[(a)]
    \item $\prod_j'\sigma_j(W_k(z^*)) = \|v_k\|\,\|w_k\|$ and
      $J_{kl} = \partial_lp_k(z^*) = w_k^TA_{k,l}v_k$: the structural content of the
      multiparameter method is $\sum_k\log(\|v_k\|\|w_k\|) - \log|\det J|$, of the same
      form as for the resultant methods (Lemma~\ref{lem:hiddenvariable}), and in
      particular local;
    \item for \textsc{MinUnif} in two variables the null vectors are explicit:
      $v_k = (1,x,\dots,x^n;\,b^{(k,y)})$ and $w_k = (b^{(k,x)};\,1,y,\dots,y^n)$, where
      $b^{(k,y)}$ collects the Horner shifts of the fiber $p_k(x,\cdot)$ deflated at the
      root, and symmetrically for $b^{(k,x)}$: the conditioned deflated-fiber objects of
      the Sylvester analysis, now with the triangular variance profile $V_j = k_{n-j}(r)$,
      for which the rank-one Schur complement generalizes verbatim:
      \[
        \E_{\mathrm{cond}}\|b\|^2 = \sum_k\frac{\bigl[\sum_{j>k}V_jr^{j-k-1}\bigr]
        \bigl[\sum_{j\le k}V_jr^j\bigr]}{\sum_jV_jr^j}\,;
      \]
    \item in the total-degree model at $d=2$,
      $\Eroots\log\|v_k\| = \frac52\log n + O(1)$, the window contributing
      $\frac32\log n$ and the collective radial excursions of
      Proposition~\ref{prop:totaldegree}(c) contributing $\log n$; hence
      \[
        \frac1N\,\E\sum_j\log\frac{\kappa_{\mathrm{MEP}}(z_j)}{\kroot(z_j)}
        = (d+1)\log n + 4\log n + O(1) = 7\log n + O(1),
      \]
      the collective-tail step being at the same sketch level as in
      Proposition~\ref{prop:totaldegree}.
  \end{enumerate}
\end{proposition}

\proofsketch
(a) At a simple singular point $\adj(W)$ has rank one with
$\|\adj(W)\|_2 = \prod_j'\sigma_j$, and Jacobi's formula
$\partial_l\det W = \tr(\adj W\,\partial_lW)$ gives the gradient identity,
$\partial_lW_k = A_{k,l}$ being constant. (b) is back-substitution in the two bidiagonal
blocks, the right-hand sides being the fiber sections $(Lv^{(1)})_i = a_{i-1}(x)$; the
Schur computation is the Sylvester one with $V_j$ in place of a constant variance. (c) On
the window the profile sums give $\E_{\mathrm{cond}}\|b\|^2\asymp n^3$; on a collective
excursion of height $t$ the products $V_jr^j$ are flat in $j$ and the conditioned sum is
$\asymp e^t(n/t)^2$, so $\log\|v\|^2\approx t + 2\log n - 2\log t$ there, and integrating
the excess against the tail $d/t$ of Proposition~\ref{prop:totaldegree}(c) adds $2\log n$
to the window value $3\log n$; the determinant and $\sigma_{\min}$ averages are those of
Proposition~\ref{prop:totaldegree}(d).
\endproof

Numerically, at total-degree roots with the pipeline of
Section~\ref{sec:verification}: the \textsc{MinUnif} determinant identity holds at machine
precision; the adjugate identity of (a) holds to $10^{-3}$ in the log at $n=11$, degrading
with the conditioning of $W_k$ at far roots but negligible against the means; and the
measured structural term and full ratio are $3.2, 3.5, 4.2, 4.7$ and $4.7, 5.7, 7.1, 8.1$
at $n=4,6,8,11$, with fitted slopes $1.5$ and $3.5$ and rising increments, the finite-size
attenuation matching that of the certified local averages of the model at equal sizes.
Two remarks are in order. First, the structural term depends on the chosen representation:
smaller or better-scaled uniform representations change the constant, and possibly the
coefficient, so the representation is a genuine design lever for the average-case
conditioning of the method. Second, the qualitative conclusion stands regardless: in the
absolute metric the multiparameter subproblem is, on average, polynomially worse
conditioned than the root, and the exponential worst cases are tail events here as
everywhere else in the family.

\subsection{Macaulay resultant: back to the overlap class}\label{sec:macaulay}

By \cite[Thm.~7.1]{graf2026numerical}, the Macaulay eigenvalue problem built with a random
linear $h$ satisfies
\[
  \keig(\lambda^*) \;\ge\;
  \frac{\|[\det Q]_\B\|_2\,\|V_{\mathrm{col}}(x^*)\|_2}{|\det J(x^*)\,h(x^*)|},
\]
where $\B$ is the basis given by the retained $h$-rows and $V_{\mathrm{col}}(x^*)$ is the
vector of column labels (monomials of degree at most $\rho := \sum_i\deg p_i - d + 1$)
evaluated at the root. Theorem~\ref{thm:reduction1} applies verbatim:
$\|[\det Q]_\B\| = |\det J|\,\|[\ell_j]_\B\|$, so
\[
  \keig(\lambda^*) \;\ge\;
  \frac{\|[\ell_j]_\B\|\,\|V_{\mathrm{col}}(x^*)\|}{|h(x^*)|}.
\]
The Macaulay method therefore sits in the overlap class together with the normal form: the
same dual norm $\|[\ell_j]_\B\| = \|V^{-T}e_j\|$, with the identities of
Lemma~\ref{lem:gram}, paired with a longer but local monomial vector
(Section~\ref{sec:outer} machinery with staircase degrees up to $\rho$) and the harmless
local factor $|h(x^*)|^{-1}$. The two-dimensional refinement of
\cite[Sec.~8]{graf2026numerical}, which bounds the numerator below by $\sigma_{\min}$ of
the Macaulay matrix, offers a complementary handle: that smallest singular value is again
controlled in expectation by Carbery--Wright, and is a natural companion statistic for the
numerical protocol.

\subsection{The hierarchy}\label{sec:hierarchy}

Table~\ref{tab:hierarchy} collects the classification established above, together with the
average-case result now available at each level; it doubles as the scoreboard of this
paper.

\begin{table}[ht]
  \centering
  \footnotesize
  \setlength{\tabcolsep}{4pt}
  \resizebox{\textwidth}{!}{%
  \begin{tabular}{lllll}
    Method & structural object & class & absolute & relative \\
    \hline
    Cayley resultant & null-vector norms, $\det J$ & local
      & $-\frac{d(d-1)}2\log n$ (Thm.~\ref{thm:cayley}) & $0$ \\
    Sylvester ($d=2$) & conditional fiber moment & local
      & $+\frac12\log n$ (\S\ref{sec:cayley}) & $+\frac12\log n$ \\
    Multiparameter EP & deflated dets of $W_k(x^*)$ & local
      & $+7\log n$ at $d=2$ (\S\ref{sec:mep}) & $+\frac{13}2\log n$ \\
    Gr\"obner elimination & coordinate pair energy & pairwise
      & $\approx-2n\log n$ at $d=2$ & $\approx-n(\log n+C)$ \\
    Rational univariate & projected pair energy & pairwise
      & $-I(u)\,N$ (Prop.~\ref{prop:rur}) & $+\delta N$ \\
    Normal form & overlap $\Theta_n(\B)$ & full config.
      & $(d{+}1)\log n{+}O(1)$, $d\le2$ (Thm.~\ref{thm:overlap})
      & $(\frac d2{+}1)\log n{+}\E\log\|M_g\|$ \\
    Macaulay resultant & overlap $\|[\ell_j]_\B\|\,\|V_{\mathrm{col}}\|$ & full config.
      & lower bounds (\S\ref{sec:macaulay}) & lower bounds
  \end{tabular}}
  \caption{Summary of the average-case results, per method: the structural object, its
    class, and the leading behaviour of the average logarithmic conditioning ratio in the
    absolute metric. Pairwise entries are stated at $d=2$; the relative column follows the
    conversion of Section~\ref{sec:relative}.}
  \label{tab:hierarchy}
\end{table}

The local terms are common to all methods, so any average-case comparison between methods
reduces to comparing structural objects. And the ordering by average-case tractability
does not follow the worst-case results of \cite{graf2026numerical}, where all methods are
exponentially unstable; the local class already contains a complete theorem
(Theorem~\ref{thm:cayley}), and the pairwise class requires only two-point information.

\subsection{Basis choices within the normal-form family}

The normal-form variants of \cite{graf2026numerical,telen2018stabilized} correspond to
different quotient bases $\B$ and different $g$. Theorem~\ref{thm:reduction1} and
Corollary~\ref{cor:wilkinson} hold for any basis, with $\Lambda_j = \Lambda_j(\B)$, and
the local terms are basis-independent. Hence the average-case comparison between variants
is the comparison of the overlap averages $\Theta_n(\B)$, and the protocol of
Section~\ref{sec:protocol} should be run with one curve per basis, on the same sampled
systems; the vertical offsets between the curves are the average-case ranking.
